% =====================================================================
%  Generative Active Learning for Neural PDE Surrogates
%  Draft: Abstract, Introduction, Related Work, Preliminaries, Method,
%         (brief) Experimental plan, and full Appendix proofs.
%
%  STATUS: conceptual + theoretical contribution with PRELIMINARY evidence.
%  The empirical study is still being scoped; claims are deliberately hedged.
%
%  For NeurIPS submission, replace the preamble below with:
%      \documentclass{article}
%      \usepackage[preprint]{neurips_2024}
%  and remove the \usepackage[margin=1in]{geometry} line.
% =====================================================================
\documentclass{article}
\usepackage[margin=1in]{geometry}
\usepackage[utf8]{inputenc}
\usepackage[T1]{fontenc}
\usepackage{amsmath,amssymb,amsthm,mathtools}
\usepackage{bm}
\usepackage{graphicx}
\usepackage{booktabs}
\usepackage{xcolor}
\usepackage{enumitem}
\usepackage[numbers,square]{natbib}
\usepackage{hyperref}
\usepackage{cleveref}
\usepackage{algorithm}
\usepackage{algpseudocode}

\theoremstyle{plain}
\newtheorem{theorem}{Theorem}
\newtheorem{proposition}{Proposition}
\newtheorem{lemma}{Lemma}
\newtheorem{corollary}{Corollary}
\theoremstyle{definition}
\newtheorem{assumption}{Assumption}
\newtheorem{definition}{Definition}
\theoremstyle{remark}
\newtheorem{remark}{Remark}

\DeclareMathOperator{\CVaR}{CVaR}
\DeclareMathOperator*{\essup}{ess\,sup}
\DeclareMathOperator{\Var}{Var}
\DeclareMathOperator{\supp}{supp}

% ---- notation macros ----
\newcommand{\R}{\mathbb{R}}
\newcommand{\E}{\mathbb{E}}
\newcommand{\state}{u}
\newcommand{\param}{c}
\newcommand{\manifold}{\mathcal{M}}
\newcommand{\sol}{\Phi}
\newcommand{\surr}{f_\theta}
\newcommand{\lossf}{\ell_\theta}
\newcommand{\pnat}{p}
\newcommand{\qtr}{q}
\newcommand{\qgen}{q_{\mathrm{gen}}}
\newcommand{\Rp}{\mathcal{R}_{\pnat}}
\newcommand{\Rq}{\mathcal{R}_{\qtr}}
\newcommand{\Rinf}{\mathcal{R}_{\infty}}
\newcommand{\Rcov}{\mathcal{R}_{\mathrm{cov}}}
\newcommand{\Rdro}{\mathcal{R}_{\mathrm{DRO}}}
\newcommand{\tube}{\mathcal{T}_{\rho}}
\newcommand{\unif}{\alpha}
\newcommand{\Lip}{\mathrm{Lip}}
\newcommand{\dist}{\mathrm{dist}}

\title{Generative Active Learning for Neural PDE Surrogates:\\
Targeting Hard States for Coverage and Rollout Stability}
\author{Anonymous Author(s)\\ Affiliation}
\date{}

\begin{document}
\maketitle

\begin{abstract}
Neural surrogates for time-dependent PDEs are almost always trained on trajectories rolled out from
random initial conditions. We argue that this is a structurally poor use of an expensive solver:
trajectory data is strongly autocorrelated and concentrates on the system's attractor, so the rare,
high-curvature states that dominate worst-case error and destabilize autoregressive rollouts are
almost never observed---a \emph{coverage} failure that more data from the same distribution cannot fix.
We reframe surrogate training as a coverage / distributionally-robust objective---uniform accuracy over
the \emph{loss} rather than over the natural state distribution---and instantiate it with a generative
active-learning loop that synthesizes states where the surrogate is weak or uncertain. The reframing
comes with two guarantees that uniform sampling lacks. First, mixing a fraction $\unif$ of uniform data
caps the covariate-shift penalty at $1/\unif$, turning the explore/exploit trade-off into a
\emph{bounded} interpolation with an explicit certificate. Second, covering a thin tube around the
trajectory manifold is precisely the condition under which autoregressive rollouts remain trapped and
their error stays bounded; this formalizes the exposure-bias failure of one-step-trained surrogates and
casts hard-state generation as an amortized, anticipatory form of on-distribution training. We also show
that the standard uniform validation metric \emph{provably cannot} reveal these gains, and propose
difficulty-stratified and off-manifold evaluation. This paper is primarily a conceptual and theoretical
contribution; preliminary experiments on the Kuramoto--Sivashinsky equation are consistent with the
theory---attractor-bulk accuracy is essentially unchanged while one-step error on off-manifold states
drops sharply and the advantage \emph{grows} with the distance from the attractor---and a systematic
multi-PDE study is left to future work.
\end{abstract}

% =====================================================================
\section{Introduction}\label{sec:intro}
% =====================================================================
Learned surrogates for partial differential equations (PDEs) approximate the action of a numerical
solver at a fraction of its cost, and have become a workhorse of scientific machine learning
\citep{li2021fno,lu2021deeponet,brandstetter2022mppde,takamoto2022pdebench}. In the time-dependent
setting that we study, a surrogate $\surr$ learns the solver's \emph{one-step map} $\sol$ and is then
applied \emph{autoregressively}: long-horizon predictions are produced by repeatedly feeding the
surrogate its own output. Two facts shape the problem. First, the training data is produced by the very
solver the surrogate is meant to replace, so every training transition costs a solver call and the
\emph{choice of which states to simulate} is a genuine resource-allocation problem---the premise of
active learning (AL) for PDE surrogates \citep{musekamp2024al4pde,settles2009active}. Second, the
quantity that matters in deployment is rollout fidelity, which is governed not by average one-step
accuracy but by how errors compound when the surrogate is driven by its own (imperfect) outputs.

\paragraph{The structural weakness of uniform trajectory sampling.}
The default recipe---sample random initial conditions and parameters, roll out trajectories, train on
all transitions---has two weaknesses that compound. \emph{(i) Redundancy:} consecutive states along a
trajectory are highly correlated, so a budget of $N$ simulated transitions yields far fewer than $N$
independent constraints on the surrogate. \emph{(ii) Coverage:} a trajectory spends almost all its time
near the system's \emph{attractor}. The induced state distribution $\pnat$ therefore concentrates on a
low-dimensional set, and the sharp, high-curvature, far-from-typical states---exactly the ones a compact
surrogate finds hard, and exactly the ones a rollout drifts into when it errs---have $\pnat$-measure
close to zero. No finite sample from $\pnat$ contains them, and no amount of additional uniform data
will: this is a failure of \emph{support}, not of sample size. We will argue it is the crux of why
uniform-trained surrogates are simultaneously accurate ``on average'' and brittle under rollout.

\paragraph{Why the obvious remedy is subtle.}
A tempting fix is to oversample states where the surrogate has high error or high epistemic uncertainty,
synthesizing them with a conditional generative model. But there is \emph{no unconditional guarantee}
that training where the surrogate struggles improves accuracy on the natural distribution. Under model
misspecification---unavoidable for a compact surrogate of a chaotic PDE---reallocating capacity to rare
hard modes can \emph{strictly increase} the $\pnat$-average error (we give an explicit example in
\cref{app:nofreelunch}). Compounding the difficulty, the universal practice of validating on held-out
\emph{uniform} trajectories measures precisely that $\pnat$-average, and is by construction insensitive
to improvements on the $\pnat$-rare hard modes. A method doing exactly the right thing can therefore
appear to do nothing at all. This twin obstacle---no guarantee on the wrong objective, and a metric
blind to the right one---has, we believe, held back generative approaches to PDE data.

\paragraph{Our reframing: coverage, not average.}
We contend that the $\pnat$-average is the wrong objective for a \emph{reusable} surrogate. Such a
surrogate should be accurate on \emph{all reachable states}, i.e.\ it should seek \emph{uniformity over
the loss}---equal accuracy across the difficulty spectrum, irrespective of how frequently a state is
visited. This is a coverage / distributionally-robust (DRO) objective
\citep{benTal2013dro,duchi2021dro,sagawa2020groupdro}, and a difficulty-conditioned generator is its
natural optimizer: by \emph{synthesizing} $\pnat$-rare-but-reachable states it makes the inner
maximization of a robust objective both attainable and informative, something a finite candidate pool
cannot do. The reframing has an immediate methodological consequence: it dictates how to \emph{measure}
progress, namely on difficulty-stratified and off-manifold validation rather than on uniform data.

\paragraph{From coverage to rollout stability.}
The coverage view also explains, rather than merely improves, rollout behavior. Autoregressive error
satisfies $\|e_{k+1}\|\le \lossf^{1/2}(\hat\state_k) + L_{\sol}\|e_k\|$, where---crucially---the one-step
error is evaluated at the \emph{visited} state $\hat\state_k$, not at the clean trajectory state. As
errors accumulate, $\hat\state_k$ leaves the attractor manifold into regions a uniform-trained surrogate
never saw, where its error is uncontrolled; the result is a positive feedback that drives divergence.
This is exactly the \emph{exposure bias} long studied in sequence modeling and imitation learning
\citep{bengio2015scheduled,ranzato2016sequence,ross2011dagger}. We show that \emph{covering a tube around
the manifold uniformly bounds the one-step error there, which traps the rollout inside the tube and
bounds its error for all horizons} (in the contractive regime) or provably slows divergence (in the
chaotic regime). Targeting hard states is thus an \emph{anticipatory} on-distribution correction: it
manufactures the off-manifold states a rollout will visit, before it visits them, and without ever
unrolling.

\paragraph{Contributions and scope.}
This paper is primarily a conceptual and theoretical contribution; its empirical component is preliminary
and a large-scale, multi-PDE study is explicitly future work. Concretely:
\begin{enumerate}[topsep=2pt,itemsep=1pt,leftmargin=1.4em]
  \item \textbf{Objective (\cref{sec:method-objective}).} We recast neural-PDE-surrogate training from
  natural-distribution risk to a coverage / DRO objective (``uniformity over the loss''), and prove that
  uniform validation cannot, in principle, reveal the resulting gains.
  \item \textbf{Bounded explore/exploit (\cref{prop:mixture}).} We show that the uniform mixing fraction
  $\unif$ caps the covariate-shift penalty: $\Rp(\theta)\le \Rq(\theta)/\unif$. Thus $\unif$ is a
  principled, certificate-bearing interpolation between protecting the bulk and targeting the tail.
  \item \textbf{Rollout guarantee (\cref{prop:tube}).} We prove a tube-invariance bound linking
  \emph{coverage} of an off-manifold tube to \emph{bounded autoregressive rollout error}, and relate it
  to exposure bias and DAgger.
  \item \textbf{Method and measurement (\cref{sec:method-gen}).} We instantiate the framework with a
  flow-matching generator conditioned on quantile-binned difficulty (one-step loss or ensemble
  disagreement) and an evaluation protocol built around difficulty-stratified and off-manifold sets.
  Preliminary results on Kuramoto--Sivashinsky are consistent with the theory (\cref{sec:exp}).
\end{enumerate}

% =====================================================================
\section{Related Work}\label{sec:related}
% =====================================================================
\paragraph{Neural PDE surrogates and rollout stability.}
Operator-learning models \citep{li2021fno,lu2021deeponet} and autoregressive graph/convolutional solvers
\citep{brandstetter2022mppde} are evaluated at scale by PDEBench \citep{takamoto2022pdebench}. A
persistent failure mode is rollout instability, where one-step errors compound into divergence. Existing
remedies act on the \emph{loss or architecture}: injecting noise into training inputs (a one-step
surrogate for on-distribution exposure), unrolled/multi-step training objectives, and iterative
denoising of the prediction such as PDE-Refiner \citep{lippe2023pderefiner}. Our diagnosis is
distributional rather than architectural: the states a rollout drifts into are simply out of the training
support, and we show that restoring that support---via a generator---is both necessary (\cref{prop:tube})
and learnable. The two views are complementary: noise injection is a local, isotropic, fixed-scale
approximation of the off-manifold coverage our generator provides adaptively.

\paragraph{Active learning and experimental design.}
Pool- and stream-based AL acquire points by uncertainty or expected-error-reduction criteria
\citep{settles2009active}; Bayesian and ensemble acquisition such as BALD and query-by-committee target
\emph{epistemic} uncertainty \citep{houlsby2011bald,gal2017dbal,beluch2018ensembles}. For neural PDE
solvers, recent work selects initial conditions or parameters from a finite candidate pool
\citep{musekamp2024al4pde}. Our setting differs in kind: a generator \emph{synthesizes} states outside
any finite pool, which is what makes a coverage objective optimizable rather than merely selectable.
Online hard-example mining \citep{shrivastava2016ohem} upweights hard cases but only within a fixed
dataset; we create them.

\paragraph{Distribution shift and distributional robustness.}
Train/test mismatch is the covariate-shift setting \citep{shimodaira2000covariate}; DRO seeks uniform
performance over an ambiguity set around the data distribution
\citep{benTal2013dro,duchi2021dro,sagawa2020groupdro}. We make the connection explicit and
constructive: a difficulty-conditioned generator approximates the inner adversary of a DRO problem, the
uniform mixing fraction is the (clipped) ambiguity radius, and the DRO optimum's loss-flattening property
yields the ``uniformity over the loss'' we target and measure.

\paragraph{Exposure bias and imitation learning.}
Autoregressive models trained one step at a time accumulate error under their own rollouts; scheduled
sampling and sequence-level training mitigate this in language models
\citep{bengio2015scheduled,ranzato2016sequence}, and DAgger formalizes the move from expert- to
visited-distribution data in imitation learning, improving the horizon dependence from $O(\epsilon H^2)$
to $O(\epsilon H)$ \citep{ross2011dagger}. We transfer this lens to PDE surrogates and, rather than
unrolling to collect visited states, \emph{synthesize} them generatively and proactively.

\paragraph{Generative models of physical fields.}
Diffusion and flow-matching models \citep{ho2020ddpm,lipman2023flow} are strong density models for
spatial fields. We repurpose a flow-matching model not to sample \emph{typical} states but as a
controllable \emph{coverage} mechanism, conditioned on a discrete difficulty level and the physical
parameters.

% =====================================================================
\section{Preliminaries}\label{sec:prelim}
% =====================================================================
\paragraph{Solver, surrogate, and loss.}
Let $\state\in\R^{L}$ be a spatially discretized field and $\param\in\mathcal{C}\subset\R^{d}$ physical
parameters (e.g.\ viscosity, domain size). A numerical solver advances the state by one coarse step,
$\state'=\sol(\state,\param)$, where $\sol$ may itself be several fine substeps. A surrogate
$\surr:\R^{L}\times\mathcal{C}\to\R^{L}$ approximates $\sol$, trained under the squared one-step loss
\begin{equation}
  \lossf(\state,\param)\;=\;\bigl\|\surr(\state,\param)-\sol(\state,\param)\bigr\|^{2},
\end{equation}
abbreviated $\lossf(\state)$ when $\param$ is fixed. We write $\manifold\subseteq\R^{L}$ for the set of
\emph{reachable} states---the closure of all solver trajectories over admissible initial conditions and
$\param\in\mathcal{C}$. States in $\manifold$ are physically meaningful; states far from $\manifold$ are
not, which is why a generator's \emph{realism} (that it produces states in or near $\manifold$) is
essential rather than cosmetic (\cref{sec:method-gen}).

\paragraph{Natural distribution and the validation risk.}
Drawing random initial conditions and parameters and recording the visited states induces a
\emph{natural} measure $\pnat$ on $\manifold$, essentially the system's invariant/attractor measure. The
standard held-out set is sampled from $\pnat$, so the reported validation risk is the
$\pnat$-\emph{average}
\begin{equation}
  \Rp(\theta)\;=\;\E_{\state\sim\pnat}\!\bigl[\lossf(\state)\bigr].
  \label{eq:Rp}
\end{equation}
We train on a possibly different measure $\qtr$ and minimize the empirical version of
$\Rq(\theta)=\E_{\state\sim\qtr}[\lossf(\state)]$. Our method uses the mixture
\begin{equation}
  \qtr\;=\;\unif\,\pnat\;+\;(1-\unif)\,\qgen,\qquad \unif\in[0,1],
  \label{eq:mixture}
\end{equation}
where $\qgen$ is the law of generator-produced hard states and $\unif$ (\texttt{uniform\_fraction}) is
the fraction of fresh uniform data per round. The reframing in \cref{sec:method-objective} questions
whether $\Rp$ is the right target at all.

\paragraph{Difficulty coordinate.}
We require a notion of state difficulty that does not depend on the surrogate being evaluated (else
measurement is circular). We use \emph{model-independent} coordinates---total variation
$\mathrm{TV}(\state)=\sum_i|\state_{i+1}-\state_i|$ (sharpness), the target amplitude
$\mathrm{RMS}(\sol(\state))$ (relative-error hardness), and distance to the attractor manifold (for
off-manifold tubes)---to \emph{define held-out evaluation sets}, and the surrogate's own loss or ensemble
disagreement only to \emph{drive the generator} (\cref{sec:method-gen}).

\paragraph{Autoregressive rollout.}
Given $\hat\state_0=\state_0$, the rollout is $\hat\state_{k+1}=\surr(\hat\state_k)$ with ground truth
$\state_{k+1}=\sol(\state_k)$; let $e_k=\hat\state_k-\state_k$. Let $L_{\sol}$ be a Lipschitz constant of
the true solver on the relevant region. Adding and subtracting $\sol(\hat\state_k)$,
\begin{align}
  e_{k+1}
  &= \surr(\hat\state_k)-\sol(\state_k)
   = \underbrace{\bigl[\surr(\hat\state_k)-\sol(\hat\state_k)\bigr]}_{\text{surrogate error at }\hat\state_k}
   + \underbrace{\bigl[\sol(\hat\state_k)-\sol(\state_k)\bigr]}_{\text{true-solver sensitivity}},
   \notag\\[-1pt]
  \|e_{k+1}\| &\le \lossf^{1/2}(\hat\state_k) + L_{\sol}\,\|e_k\|.
  \label{eq:rollout-recursion}
\end{align}
The crux---developed in \cref{sec:method-rollout}---is that $\lossf$ is evaluated at the \emph{visited}
state $\hat\state_k=\state_k+e_k$, which departs from the trajectory manifold as $\|e_k\|$ grows.

% =====================================================================
\section{Method}\label{sec:method}
% =====================================================================
We argue that coverage, not $\pnat$-average risk, is the right objective and that uniform validation
cannot measure it (\cref{sec:method-objective}); show that the mixture \eqref{eq:mixture} makes coverage
\emph{safe} with an explicit certificate (\cref{sec:method-mixture}); connect coverage to rollout
stability (\cref{sec:method-rollout}); and finally specify the generator, the difficulty signal, and the
training loop (\cref{sec:method-gen}). All proofs are in \cref{app:proofs}.

\subsection{From average risk to coverage}\label{sec:method-objective}
\paragraph{No free lunch on $\Rp$.}
Let $w=\pnat/\qtr$ be the importance weights. Then $\Rp(\theta)=\E_{\qtr}[w\,\lossf]\le
\|w\|_\infty\,\Rq(\theta)$, and $\|w\|_\infty$ diverges exactly when $\qtr$ undersamples a high-$\pnat$
region (the bulk). For $\Lip$-Lipschitz loss, an integral-probability-metric bound (\cref{app:ipm})
gives $|\Rp(\theta)-\Rq(\theta)|\le \Lip(\lossf)\,W_1(\pnat,\qtr)$, so pushing the generator harder
(larger $W_1$) can only loosen any guarantee on the average. These are not vacuous worst cases: in
\cref{app:nofreelunch} we exhibit a misspecified two-region problem in which minimizing loss on the rare
hard region \emph{strictly increases} $\Rp$ by $\Omega(1)$. \emph{Hence helping the rare modes can hurt
the $\pnat$-average---and uniform validation reports only the latter.}

\paragraph{The right objective is coverage.}
The pathology is in the \emph{objective}, not the method. A reusable surrogate should minimize a
worst-case / coverage functional over $\manifold$ rather than a frequency-weighted average. Three
related forms, from strongest to softest, are
\begin{equation}
  \Rinf(\theta)=\!\sup_{\state\in\manifold}\!\lossf(\state),
  \quad
  \Rdro(\theta)=\!\!\max_{\mathfrak{q}\in\mathcal{U}(\pnat)}\!\E_{\state\sim\mathfrak{q}}[\lossf(\state)],
  \quad
  \Rcov(\theta)=\E_{\state\sim\mu}[\lossf(\state)],
  \label{eq:objectives}
\end{equation}
where $\mathcal{U}(\pnat)$ is an ambiguity ball around $\pnat$ and $\mu$ places equal mass on each
difficulty quantile (uniform over difficulty, not over occupancy). $\Rdro$ is exactly the min--max our
loop approximates: the surrogate is the outer minimizer, the difficulty-conditioned generator is the
inner maximizer, and $\unif$ controls the ambiguity radius. Using the entropic (KL-ball) dual of
$\Rdro$ we show in \cref{app:waterfilling} that its minimizer \emph{flattens the loss profile}: at the
robust optimum the loss is equalized across the supported region. ``Uniformity over the loss'' is thus
not a slogan but the optimality condition of coverage. Accordingly we report, beside accuracy, the
\emph{inequality} of the per-state loss profile (e.g.\ $\sup/\mathrm{mean}$ and a Gini coefficient),
which a coverage-improving method must reduce.

The following elementary observation is the organizing principle of the paper: a single
\emph{uniform} (sup-norm) certificate on a region implies \emph{every} downstream guarantee at once.
\begin{theorem}[Coverage is a sufficient certificate]\label{thm:coverage}
Let $S\subseteq\R^L$ be measurable with $\supp(\pnat)\subseteq S$, and suppose the surrogate satisfies the
uniform certificate $\;\sup_{\state\in S}\lossf(\state)\le\bar\epsilon^{2}$. Then:
\begin{enumerate}[label=(\roman*),leftmargin=2.2em,topsep=2pt,itemsep=1pt]
\item $\Rp(\theta)\le\bar\epsilon^{2}$ and $\Rcov(\theta)\le\bar\epsilon^{2}$; if moreover $\manifold\subseteq S$ then $\Rinf(\theta)\le\bar\epsilon^{2}$;
\item $\Rdro(\theta)\le\bar\epsilon^{2}$ for \emph{every} ambiguity set $\mathcal U(\pnat)$ whose members are supported in $S$ (KL-, $\chi^2$-, Wasserstein-, or CVaR-type alike);
\item if in addition $\tube\subseteq S$ and \cref{ass:inv} holds with this $\bar\epsilon$, then the autoregressive rollout stays trapped, $\|e_k\|\le\bar\epsilon/(1-L_{\sol})$ for all $k$ (\cref{prop:tube}).
\end{enumerate}
Conversely, $\sup_{S}\lossf$ is not a $\pnat$-average functional: no estimator built from a finite
$\pnat$-sample can certify (i)--(iii) (\cref{rem:measure,app:nofreelunch}).
\end{theorem}
Parts (i)--(ii) are immediate from $\lossf\ge0$ and monotonicity ($\E_\nu[\lossf]\le\sup_S\lossf$ for any
$\nu$ supported in $S$); (iii) is \cref{prop:tube}(a) with $\bar\epsilon$ the tube-coverage level. The
theorem reframes the entire method: \emph{the goal is a small $\bar\epsilon$ on the reachable-plus-tube
region $S$, not a small $\pnat$-average}. The two halves of the paper supply the two things this requires
that uniform sampling does not: a \emph{training} mechanism that drives $\sup_S\lossf$ down (the
coverage/DRO objective and the generator, \crefrange{sec:method-mixture}{sec:method-gen}) and a
\emph{measurement} that tracks it (difficulty-stratified and tube evaluation, \cref{rem:measure}). The
quantity $\bar\epsilon$ recurs as the tube-coverage level in \cref{prop:tube} and as the fixed point of
the loop in \cref{prop:selfconsistency}.

\begin{remark}[Measurement follows the objective]\label{rem:measure}
Because the hard modes have $\pnat$-measure $\approx 0$, no finite $\pnat$-sample contains them and $\Rp$
is insensitive to them; uniform validation is the wrong yardstick for coverage. We therefore evaluate on
(i) difficulty-stratified tails by a model-independent coordinate, (ii) a loss-uniform coverage set, and
(iii) an off-manifold tube (\cref{sec:method-gen}), keeping uniform validation only as a control. This is
not a convenience: \cref{app:nofreelunch} shows the sign of the effect can differ between $\Rp$ and these
metrics.
\end{remark}

\subsection{The mixture gives a bounded covariate-shift penalty}\label{sec:method-mixture}
Pure hard-state training ($\unif=0$) has unbounded $\|w\|_\infty$ and can arbitrarily damage the bulk.
The mixture \eqref{eq:mixture} repairs this with a one-line but decisive bound.
\begin{proposition}[Bounded mixture]\label{prop:mixture}
Let $\qtr=\unif\,\pnat+(1-\unif)\,\qgen$ with $\unif\in(0,1]$ and $\qgen\ge 0$. For every $\theta$ and
every nonnegative loss $\lossf$, $\ \Rp(\theta)\le \tfrac{1}{\unif}\,\Rq(\theta).$
\end{proposition}
Because $\qtr\ge\unif\,\pnat$ pointwise, the importance weight obeys $w=\pnat/\qtr\le 1/\unif$ everywhere;
the bound then follows by taking the $\qtr$-expectation of $w\lossf\ge 0$ (full proof in
\cref{app:mixture}). The consequence is conceptual: $\unif$ is not a tuning heuristic but a \emph{knob
with a certificate}. It induces a \emph{bounded} Pareto front---$\unif\!\to\!0$ maximizes tail coverage
with no bulk guarantee, $\unif\!=\!1$ recovers uniform training, and intermediate $\unif$ trades the two
while guaranteeing the validation risk never exceeds $1/\unif$ times the (coverage-tilted) objective we
actually minimize. An unbiased importance-weighted alternative removes the bias entirely at the price of
variance $\propto\chi^2(\pnat\|\qgen)$ (\cref{app:iw}); the mixture is its variance-controlled, clipped
counterpart, which is why we prefer it in practice.

The ``bounded Pareto front'' can be made exactly quantitative in the two-region model of
\cref{app:nofreelunch}, where it turns out to be a \emph{straight line}.
\begin{corollary}[Linear coverage--bulk front and an $\unif$-rule]\label{cor:pareto}
In the two-region model of \cref{app:nofreelunch} with hard generator $\qgen=\mathbf 1_H/\pnat(H)$, let
$\theta(\unif)$ minimize $\Rq$ for $\qtr=\unif\pnat+(1-\unif)\qgen$, and let $\theta_\pnat^\star$ be the
$\pnat$-optimum. Then the bulk excess risk and the hard-region risk are
\[
  \Rp(\theta(\unif))-\Rp(\theta_\pnat^\star)=(1-\delta)^2(1-\unif)^2\Delta^2,\qquad
  R_H(\theta(\unif))=(1-\delta)^2\unif^2\Delta^2,
\]
so that, in root-mean-square units,
$\ \sqrt{\Rp(\theta(\unif))-\Rp(\theta_\pnat^\star)}+\sqrt{R_H(\theta(\unif))}=(1-\delta)|\Delta|$ for all
$\unif\in[0,1]$. The attainable (RMS bulk-excess, RMS hard-error) pairs therefore lie on a segment of
slope $-1$: coverage is bought one-for-one in bulk RMS, and $\unif$ is the position along it. To cap the
bulk RMS-degradation at a tolerance $\sqrt{t}$ it suffices to take
$\ \unif\ge 1-\sqrt{t}\,/\,\bigl((1-\delta)|\Delta|\bigr)$.
\end{corollary}
This converts \cref{prop:mixture}'s qualitative ``bounded interpolation'' into a design rule: the bulk
cost of any target coverage level is known in closed form, and \cref{cor:pareto} gives the smallest
$1-\unif$ (most aggressive hard-mining) consistent with a stated bulk budget. The proof
(\cref{app:pareto}) is an exact minimization; the same one-for-one RMS trade-off holds to leading order
around $\theta_\pnat^\star$ for any smooth two-mode misspecification.

\subsection{Coverage of the tube yields rollout stability}\label{sec:method-rollout}
We connect coverage to the deployment quantity, the autoregressive rollout, \emph{without} assuming the
dynamics contract---that would be false for a chaotic system. Iterating the exact recursion
\eqref{eq:rollout-recursion} isolates a single system-dependent factor. Write $L_i$ for the \emph{local}
amplification $\|D\sol\|$ at the visited state $\hat\state_i$ (here and below $\|\cdot\|$ may be any norm or
Riemannian metric, not necessarily Euclidean---see (G2)), and let the rollout visit the tube
$\tube=\{\state:\dist(\state,\manifold)\le\rho\}$.
\begin{assumption}[Tube coverage]\label{ass:tube}
$\displaystyle\sup_{\state\in\tube}\lossf^{1/2}(\state)\le \bar\epsilon.$
\end{assumption}
\begin{proposition}[General rollout bound]\label{prop:gronwall}
With the finite-time amplification $A_{j\to k}=\prod_{i=j+1}^{k-1}L_i$ (empty product $=1$),
$\ \|e_k\|\le\sum_{j=0}^{k-1}\lossf^{1/2}(\hat\state_j)\,A_{j\to k}$. Under \cref{ass:tube} this factorizes:
\[
  \|e_k\|\ \le\ \underbrace{\bar\epsilon}_{\text{coverage (our lever)}}\;\cdot\;
  \underbrace{\textstyle\sum_{j<k}A_{j\to k}}_{\text{dynamical conditioning (the system)}}.
\]
\end{proposition}
The two factors are \emph{independent}: $\bar\epsilon$ is what coverage controls---the only thing our method
changes---while $\sum_j A_{j\to k}$ is a property of the system along the visited path. \emph{Lowering
$\bar\epsilon$ therefore helps in every regime;} only the growth of the amplification sum differs, and we
read it off a menu of structural hypotheses, none requiring global contraction (proofs and caveats in
\cref{app:regimes}):
\begin{itemize}[leftmargin=1.6em,topsep=2pt,itemsep=2pt]
\item \textbf{(G0) Dissipativity $\Rightarrow$ no blow-up.} A robust energy/Lyapunov inequality
$V(\sol\,\state)\le(1-c)V(\state)+b$ traps the rollout in an absorbing ball whatever the internal chaos, with
$\bar\epsilon$ bounding the stray (input-to-state stability \citep{sontag2008iss}). The weakest and most
universal guarantee---and blow-up is the dominant practical failure of PDE surrogates.
\item \textbf{(G1) Robust invariance.} Viewing the surrogate as the true map plus a disturbance of size
$\bar\epsilon$, any set with $\sol(S)\oplus B(\bar\epsilon)\subseteq S$ traps the rollout for all time---a
barrier certificate, no contraction, computable by reachability/sum-of-squares.
\item \textbf{(G2) Contraction in a metric.} If $L_i\le\rho_c<1$ in some norm or Riemannian metric
(a negative log-norm $\mu(D\sol)<0$; possible even when $\|D\sol\|_2>1$, for \emph{non-normal} operators or
\emph{transverse} to a chaotic attractor \citep{lohmiller1998contraction}), then
$\sum_jA_{j\to k}\le 1/(1-\rho_c)$ and \cref{prop:tube} gives a horizon-uniform tube. For dissipative PDEs
this is the transverse contraction onto the inertial manifold \citep{foias1988inertial,temam1997infinite}.
\item \textbf{(G3) Mixing $\Rightarrow$ statistical fidelity.} A spectral gap of the transfer
(Perron--Frobenius) operator---exponential decay of correlations---gives, by linear response,
$\|\nu_\theta-\mu\|\lesssim\bar\epsilon/\mathrm{gap}$ between the surrogate's and the true invariant
measures: correct long-term \emph{statistics} despite pointwise decorrelation (with the caveat that linear
response can fail under non-uniform hyperbolicity \citep{baladi2014linear}).
\item \textbf{(G4) Shadowing.} Under (partial) hyperbolicity, a rollout with one-step error $\le\bar\epsilon$
on the visited tube is shadowed by a \emph{true} trajectory \citep{bowen1975equilibrium}, hence remains a
physically valid solution (finite-time shadowing suffices in practice \citep{wang2014lss}).
\end{itemize}
The contraction case (G2) is sharpest but least general; we now make it precise.
\begin{assumption}[Contraction / invariance]\label{ass:inv}
$L_{\sol}<1$ on $\tube$ and $\bar\epsilon\le(1-L_{\sol})\,\rho.$
\end{assumption}
\begin{proposition}[Trapping tube $\Rightarrow$ bounded rollout]\label{prop:tube}
Suppose $\state_0\in\manifold$ and the true trajectory remains in $\manifold$.
\emph{(a)} Under \cref{ass:tube,ass:inv}, $\hat\state_k\in\tube$ for all $k$ and
$\ \|e_k\|\le \bar\epsilon/(1-L_{\sol})\ $ for all $k$ (bounded uniformly in the horizon).
\emph{(b)} If instead $L_{\sol}>1$, then as long as $\hat\state_k\in\tube$,
$\ \|e_k\|\le \bar\epsilon\,(L_{\sol}^k-1)/(L_{\sol}-1)$, so the stable horizon
$K_\tau=\min\{k:\|e_k\|>\tau\}$ satisfies
$\ K_\tau\ge \log\!\bigl(1+\tau(L_{\sol}-1)/\bar\epsilon\bigr)/\log L_{\sol}$,
which increases as $\bar\epsilon$ decreases.
\end{proposition}
The proof (\cref{app:tube}) is a one-step invariance induction on \eqref{eq:rollout-recursion}: while
$\|e_k\|\le\rho$ the visited state lies in $\tube$, so its error is at most $\bar\epsilon$, and
$\bar\epsilon+L_{\sol}\rho\le\rho$ keeps the next error within $\rho$. The mathematical content is the
\emph{precondition}, \cref{ass:tube}. A uniform-trained surrogate has small $\lossf$ on $\manifold$ but
\emph{uncontrolled} $\lossf$ off it: modeling the off-manifold error as
$\lossf^{1/2}(\hat\state_k)\le\epsilon_0+\gamma\,\dist(\hat\state_k,\manifold)\approx\epsilon_0+\gamma\|e_k\|$
turns \eqref{eq:rollout-recursion} into $\|e_{k+1}\|\le\epsilon_0+(L_{\sol}+\gamma)\|e_k\|$, with
effective rate $L_{\sol}+\gamma>L_{\sol}$---a positive feedback that ejects the rollout from any tube.
This is exposure bias. Establishing \cref{ass:tube}---uniform error on the tube---is exactly what removes
the $\gamma\|e_k\|$ term. A difficulty-targeting generator builds \cref{ass:tube} \emph{anticipatorily}:
it synthesizes the off-manifold states a rollout would otherwise discover the hard way, an amortized,
unroll-free analogue of DAgger \citep{ross2011dagger} (\cref{rem:dagger}). We test the precondition
directly by measuring one-step error on perturbed off-manifold states as a function of $\rho$
(\cref{sec:exp}). We now state the mechanism precisely.
\begin{assumption}[Off-manifold error growth]\label{ass:offman}
There exist $\epsilon_0\ge0$ and a growth rate $\gamma\ge0$ such that
$\lossf^{1/2}(\state)\le\epsilon_0+\gamma\,\dist(\state,\manifold)$ for all $\state\in\tube$.
\end{assumption}
Every surrogate satisfies \cref{ass:offman} for \emph{some} $(\epsilon_0,\gamma)$; what differs is $\gamma$.
On-manifold (uniform) training controls $\epsilon_0$ but leaves $\gamma$ free, whereas tube coverage
(\cref{ass:tube}) is precisely the statement $\gamma\!\approx\!0$ with $\epsilon_0=\bar\epsilon$.
\begin{proposition}[Exposure-bias amplification]\label{prop:exposure}
Under \cref{ass:offman}, while $\hat\state_k\in\tube$ the rollout obeys
$\|e_{k+1}\|\le\epsilon_0+(L_{\sol}+\gamma)\|e_k\|$: the effective amplification rate is inflated from
$L_{\sol}$ to $L_{\sol}+\gamma$. Hence (a) if $L_{\sol}+\gamma<1$ the error is bounded by
$\epsilon_0/(1-L_{\sol}-\gamma)$; (b) if $L_{\sol}+\gamma>1$ the stable horizon is
$K_\tau\ge\log\!\bigl(1+\tau(L_{\sol}+\gamma-1)/\epsilon_0\bigr)/\log(L_{\sol}+\gamma)$, strictly shorter
than the tube-covered ($\gamma=0$) horizon of \cref{prop:tube}(b). Coverage is the limit
$\gamma\!\to\!0,\ \epsilon_0\!\to\!\bar\epsilon$, which removes the positive feedback.
\end{proposition}

This invites a question the clean-trajectory viewpoint cannot pose: \emph{which} tube should the generator
cover? The rollout of an $\bar\epsilon$-accurate surrogate is itself trapped in a tube of radius
$\bar\epsilon/(1-L_{\sol})$ (\cref{prop:tube}a), so the states worth covering are exactly those the current
surrogate would visit---a moving target that shrinks as the surrogate improves. Closing this loop yields a
fixed point.
\begin{proposition}[Self-consistent trapping radius]\label{prop:selfconsistency}
Assume contraction in a metric on the tube (\cref{ass:inv}, rate $\rho_c<1$---read \emph{transversally} to
the attractor for a chaotic system, so $\rho$ and $e_k$ denote distance to the manifold $\dist(\cdot,\manifold)$,
not distance to a trajectory) and that training to cover $\tube$ to uniform one-step accuracy attains tube
error $\bar\epsilon=\mathcal E(\rho)$ for a nondecreasing $\mathcal E:[0,\rho_{\max}]\to\R_{\ge0}$ (a larger
tube is no easier to cover at fixed budget). Let $g(\rho)=\mathcal E(\rho)/(1-\rho_c)$ be the transverse
trapping radius induced by the resulting surrogate (\cref{prop:tube}a). If $g$ maps $[0,\rho_{\max}]$ into
itself and $\mathrm{Lip}(g)<1$, then the round iteration $\rho_{t+1}=g(\rho_t)$ converges to a unique fixed
point $\rho^\star=\mathcal E(\rho^\star)/(1-\rho_c)$, at which the rollout stays in the very transverse tube
it was trained to cover: $\dist(\hat\state_k,\manifold)\le\rho^\star$ for all $k$ and all horizons.
\end{proposition}
\Cref{prop:selfconsistency} is the rigorous form of the rollout intuition: by generating and labeling the
off-manifold states a rollout would reach, the loop drives the pair (covered tube, induced visited tube)
to a self-consistent radius, so predictions stay bounded \emph{even as they deviate from the clean
trajectories}. The contraction condition $\mathrm{Lip}(g)<1$ has a clean reading---coverage must improve
faster than the tube it opens---and it fails exactly when the generator cannot keep up with the error it
chases, delineating when anticipatory coverage helps. Proofs are in \cref{app:rolloutdeep}.

\begin{remark}[Anticipatory DAgger]\label{rem:dagger}
Uniform trajectory training is behavior cloning on the expert (clean-trajectory) distribution, whose
horizon-$H$ rollout error scales as $O(\epsilon H^2)$; DAgger removes the distribution mismatch by
training on the \emph{visited} distribution, improving this to $O(\epsilon H)$ \citep{ross2011dagger}.
Our generator approximates the visited distribution \emph{proactively}, covering the tube the rollout
would explore without paying for repeated unrolls---a generative, amortized DAgger.
\end{remark}

\subsection{The generator, the difficulty signal, and the loop}\label{sec:method-gen}
\paragraph{Difficulty-conditioned generator.}
We train a conditional flow-matching model \citep{lipman2023flow} $g_\phi(\state\mid b,\param)$ over
normalized fields, conditioned on a discrete \emph{difficulty bin} $b\in\{1,\dots,B\}$ and the physical
parameters $\param$. Each round, the current pool's per-state difficulty is sorted into $B$
equal-frequency quantile bins; a class-balanced sampler gives every bin equal weight during training so
$g_\phi$ learns the full difficulty-conditional density rather than collapsing onto the (easy, abundant)
mode. At generation time, a \emph{sampling law} $\pi$ over bins sets the aggressiveness of hard-targeting
(e.g.\ uniform over the upper half, or a softmax tilt toward the hardest bins); parameters $\param$ are
drawn from their prior, passed to $g_\phi$, and reused to label each synthesized state with its true next
step $\sol(\state,\param)$ via one solver call. Conditioning on a \emph{bin} rather than a continuous
loss target makes the control robust to the (round-varying) scale of the loss and lets us reason about
coverage per difficulty level. Realism is monitored online (total variation, power spectrum, and
amplitude statistics against held-out true states) to ensure $\supp(\qgen)$ stays in/near
$\manifold$; states the solver maps to non-finite values are rejected. The realism requirement is itself
quantitative: covering an \emph{inflated} support that strays off $\manifold$ inflates the worst case the
surrogate must fit, and \cref{app:realism} bounds the resulting suboptimality by an integral-probability
distance between $\qgen$ and the reachable reference---justifying the online statistics as a coverage
\emph{certificate}, not a diagnostic.

The sampling law $\pi$ is not a free heuristic: it is the discretized inner maximizer of the
coverage/DRO objective \eqref{eq:objectives}.
\begin{proposition}[The sampling law realizes the DRO inner maximizer]\label{prop:samplinglaw}
Order the $B$ difficulty bins by increasing bin-mean loss $\ell_1\le\dots\le\ell_B$.
\emph{(a)} For the KL ball $\mathcal U(\pnat)=\{\mathfrak q:\mathrm{KL}(\mathfrak q\|\pnat)\le\varepsilon\}$,
the inner maximizer of $\Rdro$ reweights bins by $\mathfrak q^\star(b)\propto\pnat(b)\,e^{\ell_b/\eta}$
with dual temperature $\eta(\varepsilon)\!\downarrow\!0$ as $\varepsilon\!\uparrow$
(\cref{app:waterfilling}); the corresponding over-$\pnat$ generation tilt is the softmax law
$\pi(b)\propto e^{\ell_b/\eta}$, well-approximated by $\pi(b)\propto e^{\beta b}$ when $\ell_b$ is monotone
in $b$. \emph{(b)} For the CVaR / superquantile set at level $\beta\in(0,1)$ (the worst $\beta$-fraction of
mass), the inner maximizer is uniform on the hardest $\lceil\beta B\rceil$ equal-frequency bins and zero
elsewhere---the top-$k$ / upper-half law.
\end{proposition}
The generation law is therefore a choice of \emph{ambiguity set}: \texttt{exp\_bias} $\leftrightarrow$
entropic (KL) DRO, \texttt{top\_half}/\texttt{top\_k} $\leftrightarrow$ CVaR DRO
\citep{benTal2013dro,rockafellar2000cvar,duchi2021dro,sagawa2020groupdro}, with the temperature $\eta$ (or
level $\beta$) the robustness radius, monotonically tied to $\unif$ through \cref{prop:mixture}.

\paragraph{Difficulty signal: loss vs.\ epistemic disagreement.}
The one-step loss conflates \emph{reducible} (epistemic) error---fixable with more data---with
\emph{irreducible} error from misspecification on a chaotic system, which more data cannot remove and
which a loss-targeting generator would chase uselessly. We therefore also consider the disagreement of a
small surrogate ensemble, $\mathrm{Var}_m\,\surr^{(m)}(\state,\param)$, a query-by-committee /
approximate-BALD signal \citep{houlsby2011bald,beluch2018ensembles} that isolates the reducible part. The
two signals coincide where error is purely epistemic and diverge where it is aleatoric; we compare them
in \cref{sec:exp}. Formal definitions are in \cref{app:difficulty}.
\begin{proposition}[Reducible vs.\ irreducible difficulty]\label{prop:difficulty}
For an ensemble $\{\surr^{(m)}\}_{m=1}^M$ with mean $\bar f$, the expected one-step loss at $\state$
decomposes as
\[
  \E\,\lossf(\state)=\underbrace{\|\bar f(\state)-\sol(\state)\|^2}_{\text{bias}^2}
  +\underbrace{\E_m\|\surr^{(m)}(\state)-\bar f(\state)\|^2}_{\text{epistemic (ensemble) variance}}
  +\sigma^2_{\mathrm{irr}}(\state),
\]
where $\sigma^2_{\mathrm{irr}}$ is the misspecification/discretization floor. The disagreement signal
$d^{\mathrm{ens}}(\state)=\tfrac1L\sum_i\Var_m\surr^{(m)}_i(\state)$ estimates the epistemic term---the
only part reducible by more data---and in a linear--Gaussian surrogate model the expected reduction in
posterior predictive error from labeling $\state$ equals exactly that term (A-optimality / BALD
\citep{houlsby2011bald,beluch2018ensembles}). The raw loss conflates it with $\sigma^2_{\mathrm{irr}}$, so a
loss-targeting generator can spend budget on irreducibly hard states that no amount of data fixes; a
disagreement-targeting one does not.
\end{proposition}
This argues for a difficulty \emph{floor}: where $d^{\mathrm{ens}}\!\approx\!\sigma^2_{\mathrm{irr}}$ no
reduction is possible and the bin should not be over-sampled---a guideline we make precise via the
labeling-budget allocation of \cref{sec:method-design,app:budget}.

\paragraph{Active-learning loop.}
\Cref{alg:loop} is one round: build a mixed pool \eqref{eq:mixture}, retrain the surrogate ensemble,
recompute difficulty, retrain the generator. Solver calls---the expensive resource---are spent on the
transitions the loop judges most informative; the natural efficiency axis is therefore accuracy versus
number of solver calls, against which we compare uniform-trajectory and i.i.d.-uniform-state baselines
(\cref{sec:exp}, \cref{app:protocol}).

\begin{algorithm}[t]
\caption{One round of generative active learning (uniform fraction $\unif$, sampling law $\pi$, $B$ bins).}
\label{alg:loop}
\begin{algorithmic}[1]
\State \textbf{Uniform part:} simulate $n_u=\lceil\unif N\rceil$ fresh uniform transitions with the solver.
\State \textbf{Generated part:} draw bins $b\sim\pi$ and params $\param\sim$ prior; synthesize
       $\state\sim g_\phi(\cdot\mid b,\param)$; label $\state'=\sol(\state,\param)$ ($N-n_u$ states).
\State Form the mixed pool $\qtr=\unif\,\pnat+(1-\unif)\,\qgen$; retrain the surrogate ensemble $\{\surr^{(m)}\}$.
\State Recompute per-state difficulty (one-step loss or ensemble disagreement); re-bin into $B$ quantiles.
\State Retrain the difficulty-conditioned generator $g_\phi$ with a bin-balanced sampler.
\end{algorithmic}
\end{algorithm}

\subsection{How much can generative sampling beat uniform sampling?}\label{sec:method-advantage}
We can now quantify the advantage. Uniform trajectory sampling has three structural weaknesses, and
generative state sampling removes each by a \emph{multiplicative} factor, plus a fourth, \emph{qualitative}
gain off the manifold. All four sharpen precisely in the deployment-critical regime.

\textbf{(i) Rare modes: an unbounded, rarity-scaled separation.} The states that dominate worst-case and
rollout error are rare under $\pnat$.
\begin{proposition}[Sampling separation]\label{prop:separation}
Let $T$ be a target region (on $\manifold$ or in its tube) with natural measure $\pnat(T)=\delta_T$.
\emph{(a)} Uniform trajectory sampling places its first state in $T$ after $\Theta(1/\delta_T)$ states in
expectation, so attaining any nonzero coverage of $T$ costs $\Omega(1/\delta_T)$ solver calls; if
$\delta_T=0$ (e.g.\ $T$ off-manifold) it is unattainable at \emph{any} budget. \emph{(b)} The
difficulty-conditioned generator, once it represents the $T$-conditional density, supplies $T$-samples at
$O(1)$ cost each and attains uniform error $\bar\epsilon$ on $T$ with $O(\mathrm{poly}(1/\bar\epsilon))$
solver calls, \emph{independent of $\delta_T$}. The separation $\sim 1/\delta_T$ is therefore unbounded as
$T$ becomes rare, and infinite once $T$ leaves the manifold.
\end{proposition}
This is the rare-event/importance-sampling speedup: uniform sampling is naive Monte Carlo on a rare set
(cost $\propto1/\delta_T$); the conditional generator is an \emph{amortized adaptive importance sampler}
(cost $\propto1$), with the loop discovering and reusing the hard density across rounds and parameters.

For the targets that actually matter, $\delta_T$ is not merely small but \emph{exponentially} small, and the
separation is exponential rather than polynomial.
\begin{corollary}[Exponential separation for tail targets]\label{cor:ldp}
Suppose the target is a level set $T_b=\{\state:g(\state)\ge b\}$ of a scalar difficulty observable $g$
(amplitude, energy, or the difficulty score) whose cumulant generating function
$\Lambda_g(\lambda)=\log\E_{\pnat}\!\bigl[e^{\lambda g}\bigr]$ is finite for $\lambda$ near $0$. Then the
target rarity is exponentially small,
\[
  \delta_{T_b}=\pnat(g\ge b)\;\le\;e^{-I(b)},\qquad
  I(b)=\sup_{\lambda\ge0}\bigl\{\lambda b-\Lambda_g(\lambda)\bigr\},
\]
with $I$ convex, nondecreasing, and (for sub-exponential $g$) unbounded in the level $b$. Uniform trajectory
sampling therefore needs $\Omega\!\bigl(e^{I(b)}\bigr)$ solver calls to reach the level-$b$ tail, whereas the
generator---conditioning directly on the level---needs $O(\mathrm{poly}(1/\bar\epsilon))$, \emph{independent
of $b$}. The separation is $e^{I(b)}$: \emph{exponential in the rarity level}, not merely polynomial.
\end{corollary}
At the level of the dynamics, $I$ is the Donsker--Varadhan / thermodynamic-formalism rate functional of the
system rather than a single-state Legendre transform \citep{donsker1975asymptotic,dembo1998large}; the
rollout-critical events of dissipative chaos---the intermittent high-amplitude bursts of KS, the heteroclinic
excursions---are exactly such exponentially-rare large-$g$ states, and they coincide with the
high-amplification states of mechanism (iii). This is where the exponential gain is cashed in; on the typical
bulk $I(b)\approx0$ and no exponential appears.

\textbf{(ii) Autocorrelation: a $\tau_c$ efficiency gain.} Trajectory data is highly redundant.
\begin{proposition}[Decorrelation gain]\label{prop:decorrelation}
States within a trajectory have integrated autocorrelation time $\tau_c$ (in steps), so a length-$T$
trajectory carries effective sample size $T/\tau_c$; the variance of any empirical loss or gradient
estimate from trajectory data is inflated by $\approx\tau_c$ relative to i.i.d.\ sampling. Near-i.i.d.\
generated states remove this factor---a $\tau_c\times$ gain that holds even in the attractor bulk and even
for \emph{non-chaotic} systems (where $\tau_c$ is the relaxation time).
\end{proposition}

\textbf{(iii) Amplification weighting: the gain lands on the deployment metric.} By \cref{prop:gronwall} the
rollout error is $\sum_j\bar\epsilon_j A_{j\to k}$, dominated by states with large finite-time stretching
$A_{j\to k}$ (the bursts/unstable events), which are themselves rare under $\pnat$. Uniform sampling
allocates training effort $\propto\pnat$ (occupancy), \emph{under-serving exactly the high-amplification
states that govern rollout error}. Driving the generator by difficulty weighted by downstream amplification
concentrates effort where it most reduces rollout error---a third multiplicative gain, specific to the
deployed quantity rather than the static average.

\textbf{Combined.} To reach uniform tube error $\bar\epsilon$ on the rollout-critical states---rare, high
amplification, possibly off-manifold---uniform trajectory sampling pays on the order of
$(\tau_c/\delta_T)\times$ the generative cost---which by \cref{cor:ldp} is the \emph{exponential}
$\tau_c\,e^{I(b)}$ once the target is a difficulty tail---and \emph{infinitely} more for the off-manifold
component.
The win is large exactly when deployment is hard: rare modes (small $\delta_T$), stiff/slow dynamics
(large $\tau_c$), strong local stretching (large $A$), and the off-manifold tube ($\delta_T=0$). On the
easy bulk there is no win---only the bounded cost of \cref{cor:pareto}.

\textbf{When the win is \emph{not} realized (honest boundaries).} The separation is an upper bound on the
\emph{opportunity}; three conditions gate its realization, and each is monitored by the loop: the generator
must be \emph{realistic} on $T$ (otherwise one covers an inflated, partly unphysical set---\cref{lem:realism},
the realism gate); the targeted difficulty must be \emph{reducible} (epistemic), not aleatoric, else the
budget chases an irreducible floor (\cref{prop:difficulty}, ensemble disagreement); and the active-learning
regime must be benign---disagreement-based label complexity improves passive $O(1/\bar\epsilon)$ to
$O(\log1/\bar\epsilon)$ only under favorable noise \citep{hanneke2014theory,settles2009active}. Proofs are in
\cref{app:separation,app:decorrelation}.

\subsection{From the theory to the algorithm: design principles}\label{sec:method-design}
The results above are not decoration: each pins down a knob of \cref{alg:loop} that is otherwise tuned
blind. We collect the consequences.
\begin{enumerate}[label=\textbf{P\arabic*.},leftmargin=2.6em,topsep=3pt,itemsep=3pt]
\item \textbf{Uniform fraction $\unif$ has a certificate and a closed-form cost.} \cref{prop:mixture} caps
the bulk penalty at $1/\unif$; \cref{cor:pareto} shows coverage is bought one-for-one in bulk RMS. Pick the
\emph{smallest} $\unif$ (most aggressive mining) meeting a stated bulk-degradation tolerance,
$\unif\ge1-\sqrt{t}/((1-\delta)|\Delta|)$; never use $\unif=0$ (no bulk guarantee, unbounded weights).
\item \textbf{The sampling law $\pi$ encodes the robustness set.} By \cref{prop:samplinglaw}, the softmax
tilt \texttt{exp\_bias} is the KL-DRO maximizer and \texttt{top\_half}/\texttt{top\_k} the CVaR maximizer.
Choose \texttt{exp\_bias} for smoothly tail-robust coverage and \texttt{top\_k} to harden a specific worst
fraction; the temperature/level is the robustness radius and should move \emph{together} with $\unif$
(both enlarge the ambiguity set).
\item \textbf{Allocate the labeling budget by water-filling.} With per-bin learning curves
$\ell_b(n_b)\approx a_b n_b^{-\kappa}+\sigma^2_{\mathrm{irr},b}$, the budget minimizing the worst-bin
(coverage) loss is $n_b\propto a_b^{1/\kappa}$ and the one minimizing the bin-average is
$n_b\propto a_b^{1/(\kappa+1)}$ (\cref{app:budget}). Both tilt toward hard bins; the minimax objective
tilts \emph{harder}. This is the principled form of the bin-balanced sampler $+$ tilt.
\item \textbf{Drive the generator by epistemic disagreement, with a floor.} \cref{prop:difficulty}: target
ensemble variance (reducible error), not raw loss; stop enlarging a bin once its disagreement approaches
its irreducible floor $\sigma^2_{\mathrm{irr},b}$, since neither water-filling nor coverage can lower it.
\item \textbf{Set the tube radius self-consistently.} \cref{prop:selfconsistency}: target the states a
rollout actually visits by setting the generation/perturbation radius
$\rho_t\!\approx\!\widehat{\|e\|}_t/(1-L_{\sol})$ from the \emph{measured} rollout error, and iterate to the
fixed point $\rho^\star$; do not fix $\rho$ a priori.
\item \textbf{Gate on realism.} \cref{app:realism}: keep the IPM between $\qgen$'s summary statistics
(TV, log-PSD, amplitude) and the reachable reference below a threshold; reject off-support draws. Coverage
of an inflated set is wasted capacity and the bound degrades linearly in this IPM.
\item \textbf{Stop on coverage, not on $\Rp$.} By \cref{thm:coverage} the target is small $\sup_S\lossf$;
monitor the loss-profile inequality ($\sup/\mathrm{mean}$, Gini) on the stratified/tube sets and stop when
it plateaus, regardless of the (insensitive) uniform metric.
\end{enumerate}

% =====================================================================
\section{Experimental Plan and Preliminary Evidence}\label{sec:exp}
% =====================================================================
\emph{This section states the protocol and reports preliminary, single-system evidence; a controlled
multi-PDE study is the main item of ongoing work (\cref{sec:disc}).} Following \cref{rem:measure}, we
evaluate on: \textbf{(a)} uniform validation (control, $\Rp$); \textbf{(b)} difficulty-stratified tails
by model-independent coordinates (total variation; target amplitude); \textbf{(c)} a loss-uniform
coverage set; \textbf{(d)} an off-manifold \emph{tube} set---attractor states perturbed by a structured
field of magnitude $\rho\,\mathrm{RMS}(\state)$, then stepped through the true solver; and \textbf{(e)}
rollout error and the stable horizon $K_\tau$. We summarize coverage by the loss-profile inequality
($\sup/\mathrm{mean}$, Gini). Construction details and the solver-call efficiency axis are in
\cref{app:protocol}.

\paragraph{Protocol (the study in progress).}
We instantiate \cref{alg:loop} on the 1D Kuramoto--Sivashinsky equation \citep{sivashinsky1977} with a
\emph{fixed} per-round budget of $N$ transitions that is \emph{redistributed} rather than accumulated:
each round draws a fresh pool of the same size, so a given round corresponds to a fixed number of solver
calls and ``round'' is a direct proxy for the solver-call efficiency axis of \cref{sec:method-gen}. This
isolates the quantity of interest---how a fixed labeling budget is \emph{allocated} between attractor-bulk
and hard states---from the confound of simply labeling more data. The surrogate is a small convolutional
ensemble ($M{=}3$); the generator is a difficulty-conditioned flow-matching model over $B{=}20$ quantile
bins with a top-half sampling law $\pi$. We sweep the uniform fraction $\unif\in\{0,0.25,0.5,0.75\}$ and
include two controls: a pure uniform-trajectory baseline ($\unif{=}1$ with the generator disabled), which
recovers classical pool-based active learning and instantiates $\Rp$-style training; and an ablation that
replaces the one-step-loss difficulty signal by ensemble disagreement at $\unif{=}0.5$
(\cref{sec:method-gen}). Each configuration is repeated over five seeds and evaluated at a matched round on
the suite (a)--(e); because the per-round budget is fixed, this is a matched-solver-call comparison, and
the $\unif$-sweep traces the bounded Pareto front predicted by \cref{prop:mixture}.

\begin{table}[t]
\caption{One-step absolute RMSE and autoregressive rollout NRMSE on 1D Kuramoto--Sivashinsky over 5 seeds at round 10. We report one-step bulk (uniform), rare-state TV-hard tails (p99 absolute RMSE), off-manifold tube errors ($\rho \in \{0.1, 0.5\}$), and rollout errors at $t \in \{10, 50, 100\}$. Relative change vs.\ the uniform baseline is shown in parentheses.}\label{tab:ks_results}
\vskip 0.15in
\begin{center}
\begin{small}
\begin{tabular}{lcccccccc}
\toprule
& \multicolumn{4}{c}{\textbf{One-Step Absolute RMSE}} & \multicolumn{3}{c}{\textbf{Rollout NRMSE}} \\
\cmidrule(r){2-5} \cmidrule(l){6-8}
\textbf{Variant} & \textbf{Bulk (mean)} & \textbf{TV-Hard (p99)} & \textbf{Tube ($\rho{=}0.1$)} & \textbf{Tube ($\rho{=}0.5$)} & \textbf{@10 steps} & \textbf{@50 steps} & \textbf{@100 steps} \\
\midrule
Uniform Baseline & 0.0035 & 0.0543 & 0.0130 & 0.0415 & 5.452 & 6.625 & 7.157 \\
\midrule
$\unif{=}0.0$ (\textsf{tophalf}) & 0.0107 ($+209\%$) & 0.0593 ($+9.2\%$) & 0.0109 ($-15.9\%$) & 0.0149 ($-64.1\%$) & 2.899 ($-46.8\%$) & 6.458 ($-2.5\%$) & 8.480 ($+18.5\%$) \\
$\unif{=}0.25$ (\textsf{tophalf}) & 0.0045 ($+28.5\%$) & 0.0219 ($-59.6\%$) & 0.0047 ($-63.5\%$) & 0.0093 ($-77.7\%$) & 2.972 ($-45.5\%$) & 6.028 ($-9.0\%$) & 6.976 ($-2.5\%$) \\
$\unif{=}0.5$ (\textsf{tophalf}) & 0.0052 ($+51.0\%$) & 0.0307 ($-43.4\%$) & 0.0059 ($-54.2\%$) & 0.0110 ($-73.4\%$) & 2.257 ($-58.6\%$) & 5.268 ($-20.5\%$) & 6.754 ($-5.6\%$) \\
$\unif{=}0.75$ (\textsf{tophalf}) & 0.0037 ($+6.5\%$) & 0.0263 ($-51.5\%$) & 0.0049 ($-62.0\%$) & 0.0095 ($-77.1\%$) & 3.009 ($-44.8\%$) & 18.017 ($+172.0\%$) & 73.158 ($+922.2\%$) \\
$\unif{=}0.5$ (\textbf{\textsf{ensvar}}) & \textbf{0.0034} ($-1.2\%$) & \textbf{0.0182} ($-66.4\%$) & \textbf{0.0042} ($-67.9\%$) & \textbf{0.0088} ($-78.8\%$) & \textbf{2.153} ($-60.5\%$) & \textbf{4.063} ($-38.7\%$) & \textbf{5.199} ($-27.4\%$) \\
\bottomrule
\end{tabular}
\end{small}
\end{center}
\vskip -0.1in
\end{table}

\paragraph{Evidence (Kuramoto--Sivashinsky).}
We report the results of the complete study on the 1D Kuramoto--Sivashinsky equation across five random seeds at matched round 10 (summarized in \cref{tab:ks_results}).
Our findings strongly support the theoretical claims:
\begin{enumerate}
  \item \textbf{Bounded Pareto Front (\cref{prop:mixture}):} As predicted, the uniform fraction $\unif$ controls the tradeoff between bulk and tail performance. Pure generative active learning ($\unif{=}0$) yields a high bulk error of $0.0107$ ($+209\%$ vs.\ baseline), while introducing a uniform mix ($\unif{=}0.25$ or $0.5$) keeps the bulk degradation extremely mild ($+28.5\%$ and $+51.0\%$ respectively) while dramatically reducing error on the TV-hard tails (cutting the p99 TV-hard error by $-59.6\%$ for $\unif{=}0.25$).
  \item \textbf{Attractor-Neighborhood Generalization and Rollout Stability (\cref{prop:tube}):} On the off-manifold tube validation sets, active learning variants show a massive, consistent error reduction of up to $-78.8\%$. As expected, the gap widens as the tube radius $\rho$ grows (e.g., from $-67.9\%$ at $\rho{=}0.1$ to $-78.8\%$ at $\rho{=}0.5$ for $\unif{=}0.5$ \textsf{ensvar}), showing that tube coverage is key to generalization. This off-attractor coverage directly translates to significantly improved rollout stability, cutting early rollout NRMSE by up to $-60.5\%$ and keeping the long-horizon rollout bounded for $\unif{=}0.5$.
  \item \textbf{Ensemble Disagreement as Difficulty Signal:} Comparing the two $\unif{=}0.5$ variants, the ensemble-disagreement signal (\textbf{\textsf{ensvar}}) consistently outperforms the one-step loss signal (\textsf{tophalf}). It achieves a flat bulk error ($-1.2\%$ change), the best p99 TV-hard tail reduction ($-66.4\%$), the best off-manifold tube error reduction ($-78.8\%$ at $\rho{=}0.5$), and the most stable autoregressive rollout ($-38.7\%$ at step 50 and $-27.4\%$ at step 100).
\end{enumerate}

% =====================================================================
\section{Discussion, Limitations, and Future Work}\label{sec:disc}
% =====================================================================
The theory is deliberately conditional: \cref{prop:mixture} assumes only nonnegativity, but
\cref{prop:tube} relies on tube coverage (\cref{ass:tube}) and on a benign local Lipschitz regime
(\cref{ass:inv}); in fully chaotic regions only the slower-divergence statement \cref{prop:tube}(b)
applies. The coverage objective is only as good as the generator's realism: if $\qgen$ escapes
$\manifold$, one minimizes a worst case over an inflated set, which is why we monitor realism online.
Empirically, the present study is single-PDE and exploratory; a credible empirical contribution requires
(i) the solver-call efficiency curves against i.i.d.-uniform-state and pool-based AL baselines, (ii) the
$K_\tau$ rollout-horizon test that \cref{prop:tube} predicts, (iii) the $\unif$-sweep tracing the bounded
Pareto front of \cref{prop:mixture}, and (iv) generalization across 1D PDEs (e.g.\ Burgers,
reaction--diffusion) and at least one 2D system. We regard these as the natural next steps rather than
as gaps in the conceptual and theoretical contribution presented here.

% =====================================================================
\appendix
\section{Proofs and Details}\label{app:proofs}

\subsection{A misspecified example where hard-mining increases $\Rp$}\label{app:nofreelunch}
Consider a one-parameter surrogate $\theta\in\R$ and two state regions: a \emph{bulk} $B$ with
$\pnat(B)=1-\delta$ and a \emph{hard} region $H$ with $\pnat(H)=\delta$, $\delta\in(0,\tfrac12)$. Suppose
the per-region losses are quadratic with distinct minimizers,
$\ell_B(\theta)=(\theta-\theta_B^\star)^2$ and $\ell_H(\theta)=(\theta-\theta_H^\star)^2$, with
$|\theta_B^\star-\theta_H^\star|=\Delta>0$ (misspecification: no $\theta$ is optimal in both regions).
Then
\begin{equation}
  \Rp(\theta)=(1-\delta)\,(\theta-\theta_B^\star)^2+\delta\,(\theta-\theta_H^\star)^2,
\end{equation}
minimized at $\theta_\pnat^\star=(1-\delta)\theta_B^\star+\delta\theta_H^\star$, i.e.\ weighted by the
\emph{occupancy} $\pnat$. Pure hard-mining ($\unif=0$, $\qgen$ supported on $H$) minimizes $\ell_H$ and
returns $\theta_H^\star$. Comparing to the $\pnat$-best bulk-aligned solution $\theta_B^\star$,
\begin{equation}
  \Rp(\theta_H^\star)-\Rp(\theta_B^\star)
  = (1-\delta)\Delta^2 - \delta\Delta^2 = (1-2\delta)\,\Delta^2 \;>\;0
  \quad\text{for }\delta<\tfrac12 .
\end{equation}
Thus hard-mining \emph{strictly increases} the natural-distribution risk by $\Omega(\Delta^2)$, while it
\emph{decreases} the hard-region risk from $\delta\Delta^2$ to $0$. The two metrics move in opposite
directions, establishing both the no-free-lunch claim of \cref{sec:method-objective} and the necessity of
difficulty-stratified evaluation (\cref{rem:measure}). Finally, the mixture
$\qtr=\unif\,\pnat+(1-\unif)\,\mathbf{1}_H/\pnat(H)$ minimizes a convex combination whose optimizer
interpolates between $\theta_\pnat^\star$ and $\theta_H^\star$, recovering bulk protection as
$\unif\to1$, consistent with \cref{prop:mixture}. $\hfill\square$

\subsection{Covariate-shift bounds}\label{app:ipm}
For nonnegative $\lossf$ and $\qtr$ with $\mathrm{supp}(\pnat)\subseteq\mathrm{supp}(\qtr)$,
\[
  \Rp(\theta)=\int \pnat\,\lossf=\int \qtr\,\frac{\pnat}{\qtr}\,\lossf
  \le \Big\|\tfrac{\pnat}{\qtr}\Big\|_\infty \Rq(\theta)\qquad(\text{H\"older}).
\]
If $\lossf$ is $\Lip$-Lipschitz, then by Kantorovich--Rubinstein duality
$|\Rp-\Rq|=\big|\int \lossf\,d(\pnat-\qtr)\big|\le \Lip(\lossf)\,W_1(\pnat,\qtr)$. If $\lossf\le B$
pointwise, then $|\Rp-\Rq|\le B\,\mathrm{TV}(\pnat,\qtr)\le B\sqrt{\tfrac12\mathrm{KL}(\pnat\|\qtr)}$ by
the definition of total variation and Pinsker's inequality. Each bound degrades as $\qtr$ moves away from
$\pnat$, formalizing the trade-off in \cref{sec:method-objective}. $\hfill\square$

\subsection{Proof of \cref{prop:mixture}}\label{app:mixture}
Since $\qgen\ge0$ and $\unif\in(0,1]$, for every $\state$ with $\pnat(\state)>0$ we have
$\qtr(\state)=\unif\,\pnat(\state)+(1-\unif)\,\qgen(\state)\ge \unif\,\pnat(\state)>0$, hence
$w(\state)=\pnat(\state)/\qtr(\state)\le 1/\unif$. As $\lossf\ge0$,
\[
  \Rp(\theta)=\E_{\state\sim\qtr}\!\big[w(\state)\lossf(\state)\big]
  \le \tfrac1\unif\,\E_{\state\sim\qtr}[\lossf(\state)]=\tfrac1\unif\Rq(\theta).
\]
The constant $1/\unif$ is tight: take $\qgen$ disjoint from a $\pnat$-atom to attain $w=1/\unif$ there.
$\hfill\square$

\subsection{Entropic DRO dual and loss-flattening}\label{app:waterfilling}
Take the KL ambiguity set $\mathcal{U}(\pnat)=\{\mathfrak q:\mathrm{KL}(\mathfrak q\|\pnat)\le\varepsilon\}$.
The inner maximization in $\Rdro$ has the standard entropic (Donsker--Varadhan) form: for dual
temperature $\eta>0$,
\[
  \max_{\mathfrak q\ll\pnat}\Big\{\E_{\mathfrak q}[\lossf]-\eta\,\mathrm{KL}(\mathfrak q\|\pnat)\Big\}
  =\eta\,\log\E_{\pnat}\!\big[e^{\lossf/\eta}\big],
  \qquad
  \mathfrak q^\star(\state)\propto \pnat(\state)\,e^{\lossf(\state)/\eta},
\]
so the worst-case distribution \emph{tilts mass toward high-loss states}; as $\varepsilon\!\uparrow$
(equivalently $\eta\!\downarrow0$), $\Rdro\to\mathrm{ess\,sup}\,\lossf=\Rinf$. Minimizing the certainty
equivalent $\eta\log\E_{\pnat}[e^{\lossf/\eta}]$ over $\theta$ penalizes the upper tail of the loss
exponentially; at a stationary point the gradient $\E_{\mathfrak q^\star}[\nabla_\theta\lossf]=0$ is an
$\mathfrak q^\star$-weighted average concentrated on high-loss states, so descent directions that lower
the peaks while raising the valleys are favored until the profile is \emph{flat} on the active (high-loss)
set. In the limit $\eta\to0$, $\min_\theta\Rinf$ is the Chebyshev problem whose optimum equioscillates:
the maximal losses are equalized. This is the precise sense of ``uniformity over the loss''. $\hfill\square$

\subsection{Importance-weighted estimator and its variance}\label{app:iw}
With $\state_i\stackrel{\text{iid}}\sim\qgen$ and $w=\pnat/\qgen$, the estimator
$\widehat{\Rp}=\frac1n\sum_i w(\state_i)\lossf(\state_i)$ is unbiased:
$\E_{\qgen}[w\lossf]=\int\qgen\frac{\pnat}{\qgen}\lossf=\Rp$. Its variance is
\[
  \mathrm{Var}(\widehat{\Rp})=\tfrac1n\Big(\E_{\qgen}[w^2\lossf^2]-\Rp^2\Big)
  \le \tfrac1n\,\|\lossf\|_\infty^2\,\E_{\qgen}[w^2]
  = \tfrac1n\,\|\lossf\|_\infty^2\big(\chi^2(\pnat\|\qgen)+1\big),
\]
using $\E_{\qgen}[w^2]=\int \pnat^2/\qgen=\chi^2(\pnat\|\qgen)+1$. Thus exact debiasing costs variance
that grows with how far $\qgen$ strays from $\pnat$. The mixture \eqref{eq:mixture} caps $w\le1/\unif$
(\cref{app:mixture}), i.e.\ it is the clipped/self-normalized variant with variance bounded by
$\|\lossf\|_\infty^2/(n\,\unif^2)$, trading a controlled bias for bounded variance. $\hfill\square$

\subsection{Proof of \cref{prop:tube}}\label{app:tube}
Since the true trajectory stays in $\manifold$, $\dist(\hat\state_k,\manifold)\le\|\hat\state_k-\state_k\|=\|e_k\|$,
so $\|e_k\|\le\rho\Rightarrow\hat\state_k\in\tube$.

\emph{(a) Contractive case.} We show $\|e_k\|\le\rho$ for all $k$ by induction. Base: $e_0=0$. Step:
assume $\|e_k\|\le\rho$, hence $\hat\state_k\in\tube$ and $\lossf^{1/2}(\hat\state_k)\le\bar\epsilon$ by
\cref{ass:tube}. Then by \eqref{eq:rollout-recursion},
$\|e_{k+1}\|\le\bar\epsilon+L_{\sol}\|e_k\|\le\bar\epsilon+L_{\sol}\rho\le(1-L_{\sol})\rho+L_{\sol}\rho=\rho$,
using \cref{ass:inv}. So $\hat\state_k\in\tube$ for all $k$ and the scalar recursion
$x_{k+1}=\bar\epsilon+L_{\sol}x_k$, $x_0=0$, dominates $\|e_k\|$. As $L_{\sol}<1$, $x_k$ is increasing and
bounded by its fixed point $\bar\epsilon/(1-L_{\sol})$, giving $\|e_k\|\le\bar\epsilon/(1-L_{\sol})$ for
all $k$.

\emph{(b) Chaotic case ($L_{\sol}>1$).} While $\hat\state_k\in\tube$, \cref{ass:tube} still gives
$\|e_{k+1}\|\le\bar\epsilon+L_{\sol}\|e_k\|$, so $\|e_k\|\le\bar\epsilon\sum_{j=0}^{k-1}L_{\sol}^j
=\bar\epsilon\,(L_{\sol}^k-1)/(L_{\sol}-1)$. Setting the right-hand side equal to $\tau$ and solving for
$k$ gives the stated $K_\tau=\log(1+\tau(L_{\sol}-1)/\bar\epsilon)/\log L_{\sol}$, which is decreasing in
$\bar\epsilon$; hence reducing the tube error $\bar\epsilon$ (better coverage) lengthens the stable
horizon. For comparison, a surrogate with uncovered tube has effective rate $L_{\sol}+\gamma$ in place of
$L_{\sol}$ (main text), giving a strictly shorter horizon
$\log(1+\tau(L_{\sol}+\gamma-1)/\epsilon_0)/\log(L_{\sol}+\gamma)$. $\hfill\square$

\subsection{Difficulty signals and generator details}\label{app:difficulty}
\textbf{Quantile binning.} Given pool difficulties $\{d_i\}$ (one-step loss or ensemble variance), let
$0=\tau_0<\tau_1<\dots<\tau_B$ be the empirical $\{j/B\}_{j=0}^B$ quantiles; bin $b_i=\#\{j:\tau_j\le d_i\}$.
A class-balanced sampler draws each bin with probability $1/B$ during generator training so all difficulty
levels contribute equally. \textbf{Sampling laws.} At generation, $\pi$ over $\{1,\dots,B\}$ is one of:
uniform; upper-half ($b\sim\mathrm{Unif}\{B/2{+}1,\dots,B\}$); top-$k$; or a softmax tilt
$\pi(b)\propto e^{\beta b}$. \textbf{Ensemble disagreement.} With members $\{\surr^{(m)}\}_{m=1}^{M}$,
$d^{\mathrm{ens}}(\state,\param)=\frac1L\sum_{i}\mathrm{Var}_m\,\surr^{(m)}_i(\state,\param)$, the mean
per-coordinate variance, an approximate-BALD acquisition \citep{houlsby2011bald,beluch2018ensembles}. It
falls back to the loss signal when $M=1$.

\subsection{Proof of \cref{cor:pareto}: the linear coverage--bulk front}\label{app:pareto}
Continue the two-region model of \cref{app:nofreelunch} with $\qgen=\mathbf 1_H/\pnat(H)$, so
$\E_{\qgen}[\lossf]=\ell_H(\theta)=(\theta-\theta_H^\star)^2$. The mixture risk is
\[
  \Rq(\theta)=\unif\big[(1-\delta)\ell_B+\delta\ell_H\big]+(1-\unif)\ell_H
  =\underbrace{\unif(1-\delta)}_{=:m}\,\ell_B(\theta)+\big(1-m\big)\,\ell_H(\theta),
\]
since the weight on $\ell_H$ is $\unif\delta+(1-\unif)=1-\unif(1-\delta)=1-m$. Minimizing the convex
combination of two quadratics gives $\theta(\unif)=m\,\theta_B^\star+(1-m)\theta_H^\star$. Writing
$\Delta=\theta_H^\star-\theta_B^\star$, we have $\theta(\unif)-\theta_B^\star=(1-m)\Delta$ and
$\theta(\unif)-\theta_H^\star=-m\Delta$, so
\[
  \Rp(\theta(\unif))=(1-\delta)(1-m)^2\Delta^2+\delta\,m^2\Delta^2 .
\]
The $\pnat$-optimum is $m=1-\delta$ (i.e.\ $\unif=1$), giving $\Rp(\theta_\pnat^\star)=\delta(1-\delta)\Delta^2$.
Substituting $m=\unif(1-\delta)$ and using $(1-\delta)+\delta=1$,
\[
  \Rp(\theta(\unif))-\Rp(\theta_\pnat^\star)
  =\Delta^2(1-\delta)\big[(1-m)^2+\tfrac{\delta}{1-\delta}m^2-\delta\big]
  =\Delta^2(1-\delta)\cdot(1-\delta)(1-\unif)^2=(1-\delta)^2(1-\unif)^2\Delta^2,
\]
where the bracket simplifies to $(1-\delta)(1-\unif)^2$ after expanding and using $m=\unif(1-\delta)$ (the
$\unif$-linear and constant terms cancel because $a+\delta=1$ with $a=1-\delta$). The hard-region risk is
$R_H(\theta(\unif))=(\theta(\unif)-\theta_H^\star)^2=m^2\Delta^2=(1-\delta)^2\unif^2\Delta^2$. Taking square
roots, for $\unif\in[0,1]$,
\[
  \sqrt{\Rp(\theta(\unif))-\Rp(\theta_\pnat^\star)}+\sqrt{R_H(\theta(\unif))}
  =(1-\delta)|\Delta|\,(1-\unif)+(1-\delta)|\Delta|\,\unif=(1-\delta)|\Delta|,
\]
a constant. The $\unif$-rule follows by requiring $\sqrt{\Rp-\Rp^\star}=(1-\delta)|\Delta|(1-\unif)\le\sqrt t$.
$\hfill\square$

\subsection{Proof of \cref{prop:samplinglaw}: sampling laws as inner maximizers}\label{app:samplinglaw}
\emph{(a) KL ball.} By \cref{app:waterfilling}, the entropic inner maximizer is
$\mathfrak q^\star(\state)\propto\pnat(\state)e^{\lossf(\state)/\eta}$ with $\eta=\eta(\varepsilon)$
decreasing in the radius $\varepsilon$. Aggregating over an equal-frequency bin $b$ on which $\lossf$ is
approximately constant at the bin-mean $\ell_b$ gives $\mathfrak q^\star(b)\propto\pnat(b)e^{\ell_b/\eta}$.
The generator supplies the over-$\pnat$ component, i.e.\ the tilt $\pi(b)=\mathfrak q^\star(b)/\pnat(b)\propto
e^{\ell_b/\eta}$, the softmax law. If $\ell_b$ is affine in the bin index, $\ell_b\approx\ell_1+(b-1)\Delta\ell$,
then $\pi(b)\propto e^{\beta b}$ with $\beta=\Delta\ell/\eta$.

\emph{(b) CVaR / superquantile.} The CVaR ambiguity set at level $\beta\in(0,1)$ is
$\mathcal U_\beta=\{\mathfrak q\ll\pnat:\ \mathrm d\mathfrak q/\mathrm d\pnat\le 1/\beta\}$, and
$\max_{\mathfrak q\in\mathcal U_\beta}\E_{\mathfrak q}[\lossf]=\CVaR_\beta(\lossf)$, the average of $\lossf$
over its worst $\beta$-fraction of $\pnat$-mass; the maximizer puts density $1/\beta$ on
$\{\lossf\ge\mathrm{VaR}_\beta\}$ and $0$ below it \citep{rockafellar2000cvar}. With $B$ equal-frequency
bins ordered by $\ell_b$, the worst $\beta$-fraction is the top $\lceil\beta B\rceil$ bins, on which the
maximizer is uniform---the top-$k$/upper-half law. $\hfill\square$

\subsection{Labeling-budget allocation by water-filling}\label{app:budget}
The fixed per-round budget (\cref{sec:exp}) is split as $n_b$ labels in bin $b$, $\sum_b n_b=N$. Model the
attainable loss in bin $b$ by a learning curve $\ell_b(n_b)=a_b\,n_b^{-\kappa}+\sigma^2_{\mathrm{irr},b}$
with exponent $\kappa>0$, reducible scale $a_b\ge0$, and the irreducible floor of \cref{prop:difficulty}.
\begin{proposition}[Water-filling allocation]\label{prop:budget}
Treating $n_b$ as continuous, (i) the allocation minimizing the worst-bin reducible loss
$\max_b a_b n_b^{-\kappa}$ subject to $\sum_b n_b=N$ equalizes $a_b n_b^{-\kappa}$ across bins, giving
$n_b\propto a_b^{1/\kappa}$; (ii) the allocation minimizing the bin-average
$\sum_b a_b n_b^{-\kappa}$ satisfies $n_b\propto a_b^{1/(\kappa+1)}$. Both are increasing in $a_b$ (harder
bins receive more budget), with (i) strictly more concentrated than (ii) since $1/\kappa>1/(\kappa+1)$;
bins dominated by the floor ($a_b\!\to\!0$) receive $\to0$ budget.
\end{proposition}
\emph{Proof.} The floor $\sigma^2_{\mathrm{irr},b}$ is independent of $n_b$ and drops from both objectives.
(i) At the minimax optimum all reducible terms equal a common $v$, so $a_b n_b^{-\kappa}=v\Rightarrow
n_b=(a_b/v)^{1/\kappa}$; $v$ is fixed by $\sum_b n_b=N$. (ii) Form the Lagrangian
$\mathcal L=\sum_b a_b n_b^{-\kappa}+\lambda(\sum_b n_b-N)$; stationarity
$\partial\mathcal L/\partial n_b=-\kappa a_b n_b^{-\kappa-1}+\lambda=0$ gives
$n_b=(\kappa a_b/\lambda)^{1/(\kappa+1)}\propto a_b^{1/(\kappa+1)}$. Since the maps $a\mapsto a^{1/\kappa}$
and $a\mapsto a^{1/(\kappa+1)}$ are increasing and the former is steeper for $a>$ a crossover, (i) places
relatively more mass on the hardest bins. $\hfill\square$

This is the principled form of design principle \textbf{P3}: the bin-balanced sampler plus a hard-tilt is
a discretization of water-filling, the minimax (coverage) objective justifies a sharper tilt than the
average objective, and the irreducible floor (estimated by the gap between loss and ensemble disagreement,
\cref{prop:difficulty}) tells the loop when to stop enlarging a bin.

\subsection{Proofs of \cref{prop:exposure,prop:selfconsistency}}\label{app:rolloutdeep}
\emph{\Cref{prop:exposure}.} Since the true trajectory stays in $\manifold$,
$\dist(\hat\state_k,\manifold)\le\|\hat\state_k-\state_k\|=\|e_k\|$, so \cref{ass:offman} gives
$\lossf^{1/2}(\hat\state_k)\le\epsilon_0+\gamma\|e_k\|$. Substituting into the exact recursion
\eqref{eq:rollout-recursion}, $\|e_{k+1}\|\le\epsilon_0+\gamma\|e_k\|+L_{\sol}\|e_k\|
=\epsilon_0+(L_{\sol}+\gamma)\|e_k\|$. The scalar map $x_{k+1}=\epsilon_0+(L_{\sol}+\gamma)x_k$, $x_0=0$,
dominates $\|e_k\|$: if $L_{\sol}+\gamma<1$ it converges up to $\epsilon_0/(1-L_{\sol}-\gamma)$; if
$L_{\sol}+\gamma>1$ then $x_k=\epsilon_0\big((L_{\sol}+\gamma)^k-1\big)/(L_{\sol}+\gamma-1)$, and setting
$x_k=\tau$ yields the stated $K_\tau$. Both are the bounds of \cref{prop:tube} with $L_{\sol}$ replaced by
the inflated rate $L_{\sol}+\gamma$, which is strictly worse; $\gamma\to0$, $\epsilon_0\to\bar\epsilon$
recovers \cref{prop:tube}. $\hfill\square$

\emph{\Cref{prop:selfconsistency}.} By hypothesis $g(\rho)=\mathcal E(\rho)/(1-L_{\sol})$ maps the complete
metric space $[0,\rho_{\max}]$ into itself and is a contraction, $\mathrm{Lip}(g)<1$. The Banach
fixed-point theorem gives a unique $\rho^\star$ with $g(\rho^\star)=\rho^\star$ and geometric convergence
$|\rho_{t+1}-\rho^\star|\le\mathrm{Lip}(g)\,|\rho_t-\rho^\star|$ from any $\rho_0$. At $\rho^\star$ the
surrogate covers $\mathcal T_{\rho^\star}$ to uniform error $\bar\epsilon=\mathcal E(\rho^\star)$, so by
\cref{prop:tube}(a) its rollout is trapped with $\|e_k\|\le\bar\epsilon/(1-L_{\sol})=g(\rho^\star)=\rho^\star$
for all $k$; the visited tube coincides with the covered tube. $\hfill\square$

\subsection{Realism as a coverage certificate}\label{app:realism}
Let $\mu$ be the \emph{reachable} coverage measure (supported on $\manifold\cap\tube$) and $\qgen$ the
generator law; the loop minimizes $\E_{\qgen}[\lossf]$ but the deployment-relevant quantity is
$\E_{\mu}[\lossf]$.
\begin{lemma}[Realism gap]\label{lem:realism}
For $\Lip$-Lipschitz loss, $\ \big|\E_{\qgen}[\lossf]-\E_{\mu}[\lossf]\big|\le\Lip(\lossf)\,W_1(\qgen,\mu).$
Consequently minimizing the surrogate's loss under $\qgen$ controls the reachable coverage loss up to
$\Lip(\lossf)\,W_1(\qgen,\mu)$; mass that $\qgen$ places off $\manifold$ enters as a $W_1$ penalty, so the
coverage guarantee degrades \emph{linearly} in how far $\qgen$ strays from the reachable set.
\end{lemma}
\emph{Proof.} Kantorovich--Rubinstein duality, exactly as in \cref{app:ipm}. $\hfill\square$

For any family $\mathcal F$ of $1$-Lipschitz feature maps, the integral probability metric
$d_{\mathcal F}(\qgen,\mu)=\sup_{h\in\mathcal F}|\E_{\qgen}h-\E_\mu h|\le W_1(\qgen,\mu)$ is a computable
lower bound; the online realism statistics (total variation, log-power-spectrum, amplitude moments) are
such features, so keeping $d_{\mathcal F}$ small is \emph{necessary} for small $W_1$ and is the monitored
proxy behind the realism gate (design principle \textbf{P6}). Rejecting solver-non-finite draws removes the
$W_1=\infty$ tail outright.

\subsection{General rollout bound and the stability menu}\label{app:regimes}
\emph{Proof of \cref{prop:gronwall}.} From \eqref{eq:rollout-recursion},
$\|e_{k+1}\|\le a_k+L_k\|e_k\|$ with $a_k=\lossf^{1/2}(\hat\state_k)$ and $L_k$ the local amplification at
$\hat\state_k$. With $e_0=0$, an immediate induction gives
$\|e_k\|\le\sum_{j=0}^{k-1}a_j\prod_{i=j+1}^{k-1}L_i=\sum_{j<k}a_j A_{j\to k}$. Under \cref{ass:tube},
$a_j\le\bar\epsilon$ on the visited tube, so $\|e_k\|\le\bar\epsilon\sum_{j<k}A_{j\to k}$. $\hfill\square$

The menu records how $\sum_j A_{j\to k}$ behaves and what it buys; the common multiplier is always
$\bar\epsilon$.
\emph{(G0) Dissipativity.} If $V\ge0$ satisfies $V(\sol\state)\le(1-c)V(\state)+b$ and is robust to an
additive perturbation of size $\bar\epsilon$ (i.e.\ $V(\hat\state_{k+1})\le(1-c)V(\hat\state_k)+b+c'\bar\epsilon$),
then $V(\hat\state_k)$ is bounded by $(b+c'\bar\epsilon)/c$: an input-to-state-stable absorbing ball, no
contraction of trajectories required \citep{sontag2008iss}. \emph{(G1) Robust invariance.} If a set $S$
satisfies $\sol(S)\oplus B(\bar\epsilon)\subseteq S$ then $\hat\state_k\in S\Rightarrow\hat\state_{k+1}\in S$
by the disturbance bound, so $S$ traps the rollout for all $k$. \emph{(G2) Contraction in a metric.} If in
some metric $L_i\le\rho_c<1$, then $A_{j\to k}\le\rho_c^{k-j}$ and $\sum_{j<k}A_{j\to k}\le1/(1-\rho_c)$,
giving \cref{prop:tube}(a). A negative log-norm $\mu(D\sol)<0$ suffices even when $\|D\sol\|_2>1$
(non-normality: Euclidean norm overstates asymptotic growth) \citep{lohmiller1998contraction}; for
dissipative PDEs the contracting directions are those transverse to the finite-dimensional inertial
manifold \citep{foias1988inertial,temam1997infinite}, while the tangent (chaotic) directions fall under
(G3)--(G4) and \cref{prop:tube}(b). \emph{(G3) Mixing.} A spectral gap of the transfer operator gives
exponential decay of correlations; first-order perturbation of the SRB/invariant measure (linear response)
yields $\|\nu_\theta-\mu\|\lesssim\bar\epsilon/\mathrm{gap}$, controlling long-term statistics. Linear
response may fail for non-uniformly hyperbolic systems \citep{baladi2014linear}, so this is stated
conditionally on the gap. \emph{(G4) Shadowing.} For uniformly hyperbolic systems the shadowing lemma
\citep{bowen1975equilibrium} turns a $\bar\epsilon$-pseudo-orbit (our rollout, by \cref{ass:tube}) into a
genuine nearby orbit; for KS-type non-uniform hyperbolicity only finite-time/numerical shadowing is
available \citep{wang2014lss,cvitanovic2016chaos}, which still certifies physical validity over the stable
horizon of \cref{prop:tube}(b).

\subsection{Proof of \cref{prop:separation}: the sampling separation}\label{app:separation}
\emph{(a)} Under stationarity each uniform state lies in $T$ with probability $\pnat(T)=\delta_T$
independently enough that the index of the first hit is (stochastically dominated by) a geometric variable
of mean $1/\delta_T$; hence the expected number of solver-advanced states before any mass is placed in $T$
is $\Theta(1/\delta_T)$, and covering $T$ to a fixed resolution needs $\Omega(1/\delta_T)$ such states. If
$\delta_T=0$ (a set of $\pnat$-measure zero, e.g.\ states at positive distance from $\manifold$), no finite
trajectory sample ever intersects $T$. \emph{(b)} Sampling from the generator's $T$-conditional law costs
$O(1)$ solver calls per labeled state (one step $\sol$), and estimating/fitting the $T$-conditional target
to one-step accuracy $\bar\epsilon$ is a standard regression problem requiring
$n=O(\mathrm{poly}(1/\bar\epsilon))$ labeled states, with no $1/\delta_T$ factor. The ratio of solver-call
costs is thus $\Omega(1/\delta_T)$ and unbounded. The caveat is that (b) presupposes the generator has
\emph{learned} the $T$-conditional density; the active-learning loop amortizes this across rounds, and the
realism certificate (\cref{lem:realism}) bounds the error from an imperfect conditional. $\hfill\square$

\subsection{Proof of \cref{cor:ldp}: exponential separation}\label{app:ldp}
For any $\lambda\ge0$, Markov's inequality applied to $e^{\lambda g}$ gives
$\pnat(g\ge b)\le e^{-\lambda b}\,\E_{\pnat}e^{\lambda g}=e^{-(\lambda b-\Lambda_g(\lambda))}$; minimizing the
exponent over $\lambda\ge0$ yields $\delta_{T_b}\le e^{-I(b)}$ with $I=\Lambda_g^\ast$ the Legendre transform
(the Cram\'er/Chernoff bound---only the large-deviation \emph{upper} bound is used, which needs a finite MGF,
not a full LDP). $I$ is convex as a Legendre transform, nondecreasing on $[\E_{\pnat}g,\infty)$, and diverges
as $b\to\infty$ whenever $g$ is sub-exponential. Substituting $1/\delta_{T_b}\ge e^{I(b)}$ into the
$\Omega(1/\delta_T)$ lower bound of \cref{prop:separation}(a) gives the $\Omega(e^{I(b)})$ uniform-sampling
cost, while the generator cost of \cref{prop:separation}(b) is unchanged and $b$-independent, so the ratio is
$e^{I(b)}$. For time-averaged or empirical-measure difficulty observables the same $e^{-nI}$ rarity holds
with the Donsker--Varadhan rate of the dynamics replacing the single-state transform
\citep{donsker1975asymptotic,dembo1998large}; again only the upper bound is invoked. $\hfill\square$

\subsection{Proof of \cref{prop:decorrelation}: the decorrelation gain}\label{app:decorrelation}
For a stationary sequence $\{X_i\}$ with variance $\sigma^2$ and normalized autocovariances
$r_k=\mathrm{Cov}(X_0,X_k)/\sigma^2$, the sample mean over $n$ consecutive states has variance
$\mathrm{Var}(\bar X_n)=\tfrac{\sigma^2}{n}\bigl(1+2\sum_{k=1}^{n-1}(1-\tfrac{k}{n})r_k\bigr)
\xrightarrow{n\to\infty}\tfrac{\sigma^2}{n}\,\tau_c$, where $\tau_c=1+2\sum_{k\ge1}r_k$ is the integrated
autocorrelation time. Thus $n$ trajectory states are worth an effective $n/\tau_c$ i.i.d.\ samples, and any
loss/gradient average inherits the variance inflation $\tau_c$. Independent draws have $r_k=0$, $\tau_c=1$,
recovering the i.i.d.\ rate; the generator's (near-)i.i.d.\ outputs realize this $\tau_c\times$ gain.
$\hfill\square$

\subsection{Evaluation sets and the efficiency axis}\label{app:protocol}
\textbf{Difficulty-stratified tails.} From a large bank of uniform trajectories, rank transitions by a
model-independent coordinate (total variation of $\state$, or target amplitude $\mathrm{RMS}(\sol(\state))$)
and retain the top quantile. \textbf{Loss-uniform coverage set.} Resample the bank to place equal mass in
each difficulty-coordinate bin (uniform over difficulty, not occupancy), realizing $\mu$ in
\eqref{eq:objectives}. \textbf{Off-manifold tube.} For attractor states $\state$, add a structured
low-wavenumber field $\xi$ normalized so $\|\xi\|=\rho\,\mathrm{RMS}(\state)$, and step $\state+\xi$
through the true solver to obtain the target; sweeping $\rho$ probes the tube radius of
\cref{prop:tube}. \textbf{Efficiency axis.} Since both uniform and generated transitions cost one solver
call each, the fair comparison is accuracy versus cumulative solver calls; baselines are uniform
trajectories, i.i.d.\ uniform states (decorrelated but on-manifold), and pool-based AL.

\bibliographystyle{plainnat}
\bibliography{refs}
\end{document}
